\documentclass[10pt,a4paper]{amsart}
\usepackage{amsmath,amssymb,amsthm,graphicx}
% Retained preamble for the Stacks Project chapter {CH} (shared project preamble).
% The chapter file pinned-original-source.tex only carries \input{{preamble}} on its first line,
% so no preamble text lives inside the chapter itself.  The Stacks Project uses one single shared
% preamble for every chapter: the snapshot stages it three times (sources/stacks-chow/preamble.tex,
% sources/stacks-etale-cohomology/preamble.tex and focused-seeds/AGP/sources/stacks-snapshot-20261011/
% preamble.tex) and those copies are identical in size.  The block below is copied mechanically,
% line for line, from sources/stacks-chow/preamble.tex lines 14-18 and 154-217: the Xy-pic setup and
% the native theorem declarations together with the mathematical macros.  Nothing was added, removed
% or redefined; the wrapper class and the amsmath/amssymb/amsthm/graphicx base packages are supplied
% by the assembler.
% For commutative diagrams we use Xy-pic
% Chapter: stacks-cohomology.tex -- Cohomology of Sheaves.
\usepackage[all]{xy}

% We use 2cell for 2-commutative diagrams.
\xyoption{2cell}

% Theorem environments.
%
\theoremstyle{plain}
\newtheorem{theorem}[subsection]{Theorem}
\newtheorem{proposition}[subsection]{Proposition}
\newtheorem{lemma}[subsection]{Lemma}

\theoremstyle{definition}
\newtheorem{definition}[subsection]{Definition}
\newtheorem{example}[subsection]{Example}
\newtheorem{exercise}[subsection]{Exercise}
\newtheorem{situation}[subsection]{Situation}

\theoremstyle{remark}
\newtheorem{remark}[subsection]{Remark}
\newtheorem{remarks}[subsection]{Remarks}

\numberwithin{equation}{subsection}

% Macros
%
\def\lim{\mathop{\mathrm{lim}}\nolimits}
\def\colim{\mathop{\mathrm{colim}}\nolimits}
\def\Spec{\mathop{\mathrm{Spec}}}
\def\Hom{\mathop{\mathrm{Hom}}\nolimits}
\def\Ext{\mathop{\mathrm{Ext}}\nolimits}
\def\SheafHom{\mathop{\mathcal{H}\!\mathit{om}}\nolimits}
\def\SheafExt{\mathop{\mathcal{E}\!\mathit{xt}}\nolimits}
\def\Sch{\mathit{Sch}}
\def\Mor{\mathop{\mathrm{Mor}}\nolimits}
\def\Ob{\mathop{\mathrm{Ob}}\nolimits}
\def\Sh{\mathop{\mathit{Sh}}\nolimits}
\def\NL{\mathop{N\!L}\nolimits}
\def\CH{\mathop{\mathrm{CH}}\nolimits}
\def\proetale{{pro\text{-}\acute{e}tale}}
\def\etale{{\acute{e}tale}}
\def\QCoh{\mathit{QCoh}}
\def\Ker{\mathop{\mathrm{Ker}}}
\def\Im{\mathop{\mathrm{Im}}}
\def\Coker{\mathop{\mathrm{Coker}}}
\def\Coim{\mathop{\mathrm{Coim}}}

% Boxtimes
%
\DeclareMathSymbol{\boxtimes}{\mathbin}{AMSa}{"02}

%
% Macros for moduli stacks/spaces
%
\def\QCohstack{\mathcal{QC}\!\mathit{oh}}
\def\Cohstack{\mathcal{C}\!\mathit{oh}}
\def\Spacesstack{\mathcal{S}\!\mathit{paces}}
\def\Quotfunctor{\mathrm{Quot}}
\def\Hilbfunctor{\mathrm{Hilb}}
\def\Curvesstack{\mathcal{C}\!\mathit{urves}}
\def\Polarizedstack{\mathcal{P}\!\mathit{olarized}}
\def\Complexesstack{\mathcal{C}\!\mathit{omplexes}}
% \Pic is the operator that assigns to X its picard group, usage \Pic(X)
% \Picardstack_{X/B} denotes the Picard stack of X over B
% \Picardfunctor_{X/B} denotes the Picard functor of X over B
\def\Pic{\mathop{\mathrm{Pic}}\nolimits}
\def\Picardstack{\mathcal{P}\!\mathit{ic}}
\def\Picardfunctor{\mathrm{Pic}}
\def\Deformationcategory{\mathcal{D}\!\mathit{ef}}

% --- appended 2026-10-12 (Chief): the project own definitions of the reference/slogan/history
% display environments, copied mechanically from the shared Stacks preamble lines 8-12
% (\usepackage{verbatim} plus the three \newenvironment lines, which the chapter preamble keeps
% verbatim from the project preamble). Nothing added, removed or redefined: these environments
% carry no mathematical content and the project itself typesets them as comments. Without these
% three lines a counted statement whose own text uses them fails with Environment slogan undefined.
\usepackage{verbatim}
\newenvironment{reference}{\comment}{\endcomment}
\newenvironment{slogan}{\comment}{\endcomment}
\newenvironment{history}{\comment}{\endcomment}
\usepackage{cleveref}
\crefname{originalreference}{original reference}{original references}
\Crefname{originalreference}{Original reference}{Original references}
\makeatletter
\newlabel{categories-definition-dual}{{0FFP}{source}{}{sourcepointer-1}{}}
\newlabel{categories-definition-dual@cref}{{[originalreference][0][]0FFP}{[1][0][]source}{}{}{}}
\makeatother
\crefname{originalreference}{original reference}{original references}
\Crefname{originalreference}{Original reference}{Original references}
\makeatletter
\newlabel{example-dual}{{0FP9}{source}{}{sourcepointer-2}{}}
\newlabel{example-dual@cref}{{[originalreference][0][]0FP9}{[1][0][]source}{}{}{}}
\makeatother
\crefname{originalreference}{original reference}{original references}
\Crefname{originalreference}{Original reference}{Original references}
\makeatletter
\newlabel{categories-remark-left-dual-adjoint}{{0FFR}{source}{}{sourcepointer-3}{}}
\newlabel{categories-remark-left-dual-adjoint@cref}{{[originalreference][0][]0FFR}{[1][0][]source}{}{}{}}
\makeatother
\crefname{originalreference}{original reference}{original references}
\Crefname{originalreference}{Original reference}{Original references}
\makeatletter
\newlabel{modules-lemma-left-dual-module}{{0FNW}{source}{}{sourcepointer-4}{}}
\newlabel{modules-lemma-left-dual-module@cref}{{[originalreference][0][]0FNW}{[1][0][]source}{}{}{}}
\makeatother
\title{Left duals in the monoidal category of complexes are the termwise duals}
\author{The Stacks Project Authors}
\date{stacks-project:cohomology.tex:local-snapshot-20261011; GNU FDL 1.2 or later}
\begin{document}
\maketitle
\noindent\textit{Editorial note.} All delivered mathematical text is the author's own, copied verbatim from the pinned Stacks Project chapter file cohomology.tex; no statement or proof step was rewritten, reordered or removed.\par\medskip
\noindent\textit{Editorial note.} The counted span cites the lemmas recorded in assembly.external\_references with canonical tags as navigation pointers.\par\medskip
\noindent\textit{Allowed standard prerequisites (editorial).}\begin{itemize}
\item Uniqueness of left duals. Categories, Remark categories-remark-left-dual-adjoint. Left duals are unique up to unique isomorphism.
\item Module-level left duals. Modules, Lemma modules-lemma-left-dual-module. F\textasciicircum{}n is locally a direct summand of a finite free module.
\end{itemize}\medskip
\ifdefined\hypertarget\hypertarget{sourcepointer-1}{}\fi
\noindent\textit{Source reference.} Original reference 0FFP: sources/stacks-chow/tags.txt, source lines 19294--19294. Cross-chapter standard prerequisite cited by the author inside the counted span; it is used as a pointer and its own proof is not part of this item.\par\medskip
\ifdefined\hypertarget\hypertarget{sourcepointer-2}{}\fi
\noindent\textit{Source reference.} Original reference 0FP9: sources/stacks-chow/tags.txt, source lines 19594--19594. Cross-chapter standard prerequisite cited by the author inside the counted span; it is used as a pointer and its own proof is not part of this item.\par\medskip
\ifdefined\hypertarget\hypertarget{sourcepointer-3}{}\fi
\noindent\textit{Source reference.} Original reference 0FFR: sources/stacks-chow/tags.txt, source lines 19296--19296. Cross-chapter standard prerequisite cited by the author inside the counted span; it is used as a pointer and its own proof is not part of this item.\par\medskip
\ifdefined\hypertarget\hypertarget{sourcepointer-4}{}\fi
\noindent\textit{Source reference.} Original reference 0FNW: sources/stacks-chow/tags.txt, source lines 19581--19581. Cross-chapter standard prerequisite cited by the author inside the counted span; it is used as a pointer and its own proof is not part of this item.\par\medskip
\begin{lemma}
\label{lemma-left-dual-complex}
Let $(X, \mathcal{O}_X)$ be a ringed space. Let $\mathcal{F}^\bullet$
be a complex of $\mathcal{O}_X$-modules. If $\mathcal{F}^\bullet$
has a left dual in the monoidal category of complexes of
$\mathcal{O}_X$-modules
(Categories, Definition \ref{categories-definition-dual})
then $\mathcal{F}^\bullet$ is a locally bounded complex whose terms are
locally direct summands of finite free $\mathcal{O}_X$-modules
and the left dual is as constructed in Example \ref{example-dual}.
\end{lemma}

\begin{proof}
By uniqueness of left duals
(Categories, Remark \ref{categories-remark-left-dual-adjoint})
we get the final statement provided we show that $\mathcal{F}^\bullet$
is as stated. Let $\mathcal{G}^\bullet, \eta, \epsilon$ be a left dual.
Write $\eta = \sum \eta_n$ and $\epsilon = \sum \epsilon_n$
where $\eta_n : \mathcal{O}_X \to
\mathcal{F}^n \otimes_{\mathcal{O}_X} \mathcal{G}^{-n}$
and
$\epsilon_n : \mathcal{G}^{-n} \otimes_{\mathcal{O}_X} \mathcal{F}^n
\to \mathcal{O}_X$. Since
$(1 \otimes \epsilon) \circ (\eta \otimes 1) = \text{id}_{\mathcal{F}^\bullet}$
and
$(\epsilon \otimes 1) \circ (1 \otimes \eta) = \text{id}_{\mathcal{G}^\bullet}$
by Categories, Definition \ref{categories-definition-dual} we see immediately
that we have
$(1 \otimes \epsilon_n) \circ (\eta_n \otimes 1) = \text{id}_{\mathcal{F}^n}$
and
$(\epsilon_n \otimes 1) \circ (1 \otimes \eta_n) =
\text{id}_{\mathcal{G}^{-n}}$.
Hence we see that $\mathcal{F}^n$ is locally a direct summand of a finite
free $\mathcal{O}_X$-module by
Modules, Lemma \ref{modules-lemma-left-dual-module}.
Since the sum $\eta = \sum \eta_n$ is locally finite, we conclude that
$\mathcal{F}^\bullet$ is locally bounded.
\end{proof}

\end{document}
